\section{Introduction}
\label{sec:intro}

Large matrix multiplication is a recurring bottleneck in verifiable
computation.
A prover may wish to convince a verifier that committed matrices
$\mat{A},\mat{B},\mat{C}$ satisfy $\mat{A}\mat{B}=\mat{C}$ without forcing the
verifier, or an underlying SNARK circuit, to recompute the full product.
This problem appears across outsourced linear algebra, proof-carrying data
pipelines, and verifiable inference for modern neural networks, where projection
and attention layers routinely contain matrix products with dimensions
$k \ge 1{,}000$.
Na\"ive SNARK arithmetization of a single $k \times k$ matrix multiplication
generates $\mathcal{O}(k^3)$ arithmetic constraints, a cost that is impractical
even for moderate~$k$.

\paragraph{The $\mathcal{O}(k^2)$ barrier}
Freivalds' classical randomized check~\cite{freivalds1977probabilistic}
reduces the problem from $\mathcal{O}(k^3)$ to $\mathcal{O}(k^2)$: given
matrices $\mat{A}, \mat{B}, \mat{C}$ and a random vector $\vect{r}$,
one verifies $\vect{r}\mat{A}\mat{B} = \vect{r}\mat{C}$ instead of
recomputing the full product.
Several recent systems, including verifiable ML pipelines such as
Mystique~\cite{weng2021mystique}, zkML~\cite{chen2024zkml}, and the
Freivalds-based baseline in~\cite{ghodsi2017safetynets}, adopt this strategy,
encoding the vector-matrix products directly as SNARK constraints.
The resulting $\mathcal{O}(k^2)$ circuit is an improvement, yet it remains
expensive: in our optimized baseline, Freivalds already requires
	$50.5$ million R1CS constraints and $405.04\,\text{s}$ of proving time at
	$k = 4{,}096$.
Thus the gap between the $\mathcal{O}(k^2)$ SNARK cost and practical
large-matrix verification remains wide.

\paragraph{Our insight: from computing to checking}
We observe that the $\mathcal{O}(k^2)$ barrier is \emph{not} inherent to
the verification problem---it arises because existing approaches
\emph{compute} the vector-matrix products inside the SNARK circuit.
Our key idea is to shift the paradigm from ``computing in the circuit'' to
``\emph{checking} in the circuit'': the prover performs the expensive
matrix arithmetic \emph{outside} the SNARK and commits to all intermediate
results; the verifier then checks a small number of lightweight algebraic
relations that suffice to detect any inconsistency.

The technical engine behind this shift is \emph{proximity testing over
linear error-correcting codes}.
By encoding each matrix row-wise with a systematic linear code $\mathcal{C}$
of relative distance $\delta$, the vector-matrix product
$\vect{x} = \vect{r}\mat{A}$ can be verified \emph{column by column} in
the encoded domain: the linearity of $\mathcal{C}$ ensures that a single
inner product (``fold'') of one codeword column suffices to check one
coordinate of the encoded product.
If the prover cheats on any relation, the committed codewords disagree
in at least a $\delta$-fraction of positions, and a random query detects
the inconsistency with probability $\ge \delta$.
Crucially, the SNARK circuit encodes the \emph{local algebraic checks}---fold
computations, sampled RS evaluations, wire-handle openings, and
message-opening checks at $t$ queried positions---while Merkle path
verification and CP-link consistency over compact leaf handles are checked by
auxiliary verifiers bound to the same transcript.
The resulting in-circuit checking size is
$\mathcal{O}(tk+E_{\mathsf{msg}}(k))$, with
$\mathcal{O}(t\log n+E_{\mathsf{link}}(k))$ auxiliary opening/linking work.
It is \emph{linear} in the matrix dimension when $t$ is a small
security-parameter-dependent constant and no optional full-codeword backend is
used.
This is a checking-layer result: the input commitments and the
intermediate message/linking backend are explicit components of the statement,
and an end-to-end verifiable-computation system must instantiate the raw input
certification path.

\paragraph{Contributions}
We present $\protocol$, a matrix-multiplication checking layer for verifiable
computation that realizes this approach.
Our contributions are as follows.

\begin{itemize}
    \item \textbf{A novel code-based verification protocol.}
    We design a three-round public-coin protocol that combines a structured
    variant of Freivalds' randomized reduction with proximity testing over
    linear codes and Merkle-tree commitments.
	    In our CP-linked instantiation, message commitments to
	    $\vect{x},\vect{y},\vect{z}$ are fixed before the query set is derived,
	    and the circuit opens public wire handles for the sampled values used in
	    the fold checks; auxiliary Merkle and CP-link verifiers connect those
	    handles to authenticated external leaves.
	    The resulting online SNARK relation has $\mathcal{O}(tk)$ sampled
	    algebraic work for fixed security parameter, a significant reduction
	    from the $\mathcal{O}(k^2)$ of Freivalds-based approaches.
	    The construction should not be read as an end-to-end
	    commitment-generation system: certified input commitments to
	    $\mat{A},\mat{B},\mat{C}$ remain a deployment interface.
    The protocol is made non-interactive via the Fiat--Shamir
    heuristic~\cite{fiat1987prove}.

    \item \textbf{Rigorous security analysis.}
    We prove that $\protocol$ achieves perfect completeness and computational
    soundness with error bounded by
    $\epsilon_{\mathsf{SNARK}} + \epsilon_{\mathsf{cc}}
    + \epsilon_{\mathsf{bind}} + (k\!-\!1)/|\mathbb{F}|
    + 4(1\!-\!\delta)^{t}$,
    where $\delta$ is the code's relative distance,
    $\epsilon_{\mathsf{SNARK}}$ is the SNARK soundness error,
	    $\epsilon_{\mathsf{cc}}$ is the message-binding and sampled handle-linking error,
	    and $\epsilon_{\mathsf{bind}}$ is the Merkle commitment's binding error.
	    The Fiat--Shamir non-interactive variant adds the standard
	    random-oracle loss $\epsilon_{\mathsf{FS}}(q_H)$.
	    The theorem is stated relative to certified public input commitments to
	    row-wise encoded matrices and to the CP-linked intermediate backend,
    matching the committed-input viewpoint used in code-based proof systems
    such as Brakedown~\cite{golovnev2023brakedown} and
    Orion~\cite{xie2022orion}.
    We also give a concrete security-budget table that separates the reported
    scaling rows from a strict 128-bit deployment target.
    The analysis decomposes the adversary's strategy into SNARK failure,
    message/linking failure, commitment forgery, algebraic evasion, and proximity
    evasion, and shows each is negligible for standard parameter choices
    (Section~\ref{sec:security}).

    \item \textbf{Concrete efficiency gains.}
	    We implement the multiplication-checking layer of $\protocol$ in Rust
	    and evaluate it against a direct Freivalds SNARK baseline across matrix
	    dimensions up to $k=2^{13}$ with the CP-linked sampled-opening backend.
    Key findings (Section~\ref{sec:implementation}):
    \begin{itemize}[nosep]
        \item \emph{Constraint reduction.}
	        The implemented checking-layer constraints grow as
	        $\mathcal{O}(k)$ for fixed $t$, versus
	        $\mathcal{O}(k^2)$ for Freivalds.
        The gap reaches $7.5\times$ at $k = 1{,}024$
        ($419{,}139$ vs.\ $3{,}156{,}298$ constraints) and
        $30.3\times$ at $k = 4{,}096$
        ($1{,}667{,}305$ vs.\ $50{,}495{,}818$ constraints);
	        the checking-layer implementation also scales to $k = 8{,}192$ with
	        $3{,}331{,}779$ constraints.

        \item \emph{Proving-time speedup.}
        The current QA-NIZK implementation $\fmQA$ outperforms Freivalds from
        $k = 256$ onward.
        At $k = 4{,}096$, prove times are $56.44\,\text{s}$ for $\fmQA$ and
        $405.04\,\text{s}$ for Freivalds, a $7.18\times$ speedup.
        Online proof bytes, excluding setup/key/CRS material, are $104$\,KB at
        $k=4{,}096$ and $126$\,KB at $k=8{,}192$.

        \item \emph{GPT-2 workload case study.}
        For a GPT-2 workload with $36$ matrix-multiplication claims at sequence
        length $1{,}024$, $\fmQA$ reduces constraints by $2.30\times$ and
        proving time by $1.77\times$ relative to Freivalds, while verification
        is slower because the current QA-NIZK implementation checks sampled
        opening and linking material.

        \item \emph{Negligible encoding overhead.}
        The prover's off-circuit encoding step takes less than
        $1\,\text{s}$ even at $k = 2^{20}$, confirming that encoding
        does not introduce a new bottleneck.
    \end{itemize}
\end{itemize}

\paragraph{Paper organization}
Section~\ref{sec:tech-overview} gives a high-level technical overview.
Section~\ref{sec:prelim} recalls the necessary background.
Section~\ref{sec:construction} presents the formal construction and complexity
analysis.
Section~\ref{sec:security} provides the security proof.
Section~\ref{sec:implementation} reports our implementation and evaluation.
Section~\ref{sec:conclusion} concludes with related-work positioning,
limitations, and future directions.
